% Transcription — fonds Alexandre Grothendieck, shelfmark 23, batch 1 (pages 1-20).
% Established from the facsimile archives/batches/23/batch-01.pdf (a mirror of
% https://grothendieck.umontpellier.fr/23.pdf), per the skill
% transcribe-grothendieck. The batch file carries no cover sheet: PDF page N is
% archivists' page N. Checked against the pencilled numbers bottom-left — "9",
% "10", "11", "15", "18", "19" and "20" read plainly on the manuscript
% leaves, at the PDF page of the same number — the alignment holds.
%
% Pass: Fable 5.1 (claude-fable-5-1), 2026-09-21 — first pass, unchecked against the pages by a human.
%
% What the batch is. Twenty leaves of manuscript notes in ink, all in his
% hand, sorted by him into a folder with sub-folders: four otherwise blank
% leaves carry only a title, top right, and cover the run that follows.
% They get no \page{}; their words become the section headings, each with
% a \note saying which leaf it comes off.
%
%   1      Cover of the whole shelfmark: « EGA IV, Calcul diff.
%          (compléments) ». No \page.
%   2      Cover « Discriminant etc ». No \page.
%   3-5    Three leaves on the canonical bundle of a complete intersection
%          X in a projective bundle P(E) over S: the exact sequences for
%          Ω¹_{P/S} and J/J², the isomorphism Ω^n_{X/S} ≃ f^*(Λ^{r+1}E ⊗
%          ⊗ Λ^{α_i}M̌_{d_i})(Σα_i d_i − r − 1), the case Σα_i d_i = r+1,
%          and the worked example of a plane cubic (E of rank 3, M ≃ L^⊗6).
%          Page 5 is a rougher draft on the discriminant δ_{r,k} of a
%          hypersurface of degree k in P^r, as a section of O((k−1)(r+1)).
%   6      Blank. Absent.
%   7      Cover « Exponentielle ». No \page.
%   8-16   « Sorites sur l'exponentielle », numbered (1) to (4) by him
%          with a (6) on page 12 — no (5) was found. exp and log between
%          topologically nilpotent elements and 1 + N(U) in a Q-algebra;
%          then in a bigebra U, primitive elements π(U) against
%          group-like elements γ(U) (Δ(g) = g ⊗ g), the bijection
%          π ⇄ γ, the transport of the group law to π by
%          Campbell-Hausdorff; then U ⊗_k Λ for Λ an infinitesimal
%          k-algebra with augmentation ideal I(Λ), and the bijection
%          π(U) ⊗ I(Λ) ⇄ 1 + π_0(U ⊗ Λ). Page 13 specialises to
%          U = the enveloping algebra of L_k(M) and reads off:
%          ξ a derivation iff exp ξ an automorphism. Pages 14-16 state
%          the Proposition: for k ⊃ Q, the Λ-automorphisms of A ⊗_k Λ
%          inducing the identity mod I(Λ) correspond to k-derivations
%          A → A ⊗_k I(Λ), with the group law of Hausdorff-Campbell;
%          when Λ is k-free of finite rank this is
%          T_{A/k} ⊗_k I(Λ). A long marginal note on page 12 sketches
%          the theorem of Cartier (U ≃ U(g) for g = π(U)) and
%          Poincaré-Witt.
%   17     Cover « Fibrés des germes d'ordre n de sections, etc.
%          Transformations infinitésimales ». No \page.
%   18-20  S a prescheme, Λ a locally free sheaf of augmented O_S-algebras
%          « de type infinitésimal » (I^n = 0), T = Spec Λ with its
%          section S → T, X a T-scheme, X_0 = X ×_T S, ξ = X_0 ⊗ Λ the
%          trivial deformation. The sheaf Aut^{X_0}_T(X) of
%          T-automorphisms inducing the identity on X_0, the scheme
%          X̄ = Hom_{T/S}(T, X) of sections, the projection X̄ → X_0,
%          the sheaf S(X̄/X_0) of sections, and the Proposition (page 20)
%          that Is^{X_0}_T(ξ, X) → S(X̄/X_0), u ↦ u(s_0), is an
%          isomorphism, with two corollaries. Page 20 breaks off
%          mid-sentence; the run continues into batch 2.
%
% Notation fixed for the batch. His underlined letters — O_S, O_X, Λ, I,
% Aut, Is, S, Hom — are rendered \underline{...}; the horizontal bar
% through Φ (page 5) and through I (page 11) is rendered the same way and
% noted at first occurrence. The cursive looped letter for the tangent
% sheaf (pages 4, 15, 16) is \mathcal{T}: page 15 explicitates it as
% Hom_B(Ω¹_{B/Λ}, B), which settles the reading. Group-like elements are
% γ(U) and primitive elements π(U) throughout, as pages 10-12 write them;
% the gothic letter of the margin of page 12 is \mathfrak{g}. The
% « Symm », « Sym », « Hom », « Aut », « Is », « gr », « Spec » are
% \operatorname / \mathrm as he writes them, without regularising Symm to
% Sym. The exp/log pair is always \underset{\log}{\overset{\exp}
% {\rightleftarrows}}.
%
% Legibility. Pages 3, 4, 8, 14-16 and 18-20 are written fair and read
% nearly end to end. Pages 5, 9-13 are a fast draft, heavily struck and
% overwritten, whose formulas carry and whose connective prose largely
% does not; the apparatus there is dense, and honestly so. The margin of
% page 12 (written bottom to top over the full height) and the corner of
% page 20 (five slanted lines around a proper name read « Greenberg »)
% are the two places a checker with the leaf in hand will gain most.
%
% What is NOT here. No attempt was made to identify which paragraph of
% EGA IV these notes prepared, nor to date them beyond the inventory's
% « 1966-1968 »; no date appears on the twenty leaves. Page 6 (blank)
% and the four cover leaves are the only pages absent from the numbering.
\documentclass[11pt,a4paper]{article}
% Le préambule commun à toutes les transcriptions du fonds Grothendieck.
%
% Il tient en une idée : **l'incertitude se compose, elle ne se résout pas.**
% Une transcription qui lisse un mot douteux a détruit la seule chose pour
% laquelle on la faisait — un lecteur doit pouvoir savoir, à chaque endroit,
% ce qui était sur le feuillet et ce qui a été ajouté. Les six macros qui
% suivent sont donc l'appareil critique, et non de la décoration.
%
% Les mêmes commandes sont reconnues par `scripts/render.mjs`, qui produit la
% vue de lecture du site. Une macro ajoutée ici doit l'être aussi là-bas,
% faute de quoi le rendu échouera bruyamment — ce qui est voulu : un rendu
% silencieusement faux est pire qu'un rendu absent.

\ProvidesPackage{grothendieck}[2026/08/08 Transcriptions du fonds Grothendieck]

\usepackage{amsmath,amssymb,amsthm}
\usepackage{mathrsfs}
\usepackage{stmaryrd}
% Les diagrammes commutatifs sont partout dans ce fonds : c'est la langue
% même de la géométrie algébrique de Grothendieck, et une transcription qui
% les rendrait en prose ne transcrirait rien.
\usepackage{tikz-cd}
% Français par défaut — c'est la langue des feuillets — mais l'anglais est
% chargé aussi : la traduction et le résumé sont des documents anglais, et la
% typographie française (espaces avant les deux-points) y serait une faute.
% Ces documents basculent par \selectlanguage{english} en tête de corps.
\usepackage[english,french]{babel}
\usepackage{fontspec}
\usepackage{geometry}
\geometry{a4paper,margin=25mm,marginparwidth=28mm}
\usepackage{marginnote}
\usepackage{xcolor}
% ulem, not soul. `soul`'s \st and \ul are built on a character-by-character
% scan that XeLaTeX's font machinery breaks: the document fails with an
% undefined control sequence rather than degrading. ulem does the same two jobs
% and is Unicode-engine safe. `normalem` stops it hijacking \emph, which the
% transcriptions use for Grothendieck's own emphasis.
\usepackage[normalem]{ulem}
\usepackage{hyperref}
\hypersetup{colorlinks=true,linkcolor=black,urlcolor=black,pdfborder={0 0 0}}

% --- Métadonnées du lot -----------------------------------------------------
% Reprises telles quelles par le script de rendu, qui les affiche en tête de
% la vue de lecture. Elles fixent ce que le fichier prétend couvrir : sans
% elles, une transcription flotte sans point d'attache dans un fonds de seize
% mille feuillets.

\newcommand{\folder}[1]{\gdef\grothFolder{#1}}
\newcommand{\batch}[1]{\gdef\grothBatch{#1}}
\newcommand{\leaves}[2]{\gdef\grothFirst{#1}\gdef\grothLast{#2}}
\newcommand{\foldertitle}[1]{\gdef\grothFoldertitle{#1}}
\gdef\grothFolder{?}\gdef\grothBatch{?}\gdef\grothFirst{?}\gdef\grothLast{?}\gdef\grothFoldertitle{}

% --- Le résumé --------------------------------------------------------------
%
% Il ouvre la lecture modernisée, sous le titre « Résumé », et il est composé
% en petit. Ce n'est pas un ornement : c'est le seul passage du document qui
% ne soit pas des mathématiques, et un lecteur doit pouvoir le reconnaître
% avant de l'avoir lu — puis le sauter s'il connaît déjà le sujet, ou s'y
% arrêter s'il ne le connaît pas. Le filet qui le suit marque la reprise du
% corps.

\newenvironment{resume}
  {\small\setlength{\parskip}{0.45em}}
  {\par\medskip\hrule\medskip}

% --- Mots-clés ---------------------------------------------------------------
%
% En anglais, en clôture du résumé de la lecture modernisée : le vocabulaire
% moderne sous lequel chercher ce que les feuillets construisent. Le manifeste
% du site (`npm run manifest`) les extrait pour étiqueter la cote entière —
% cette ligne est la source unique des tags, il n'y a pas de fichier à côté.
\newcommand{\keywords}[1]{\par\smallskip\noindent
  {\footnotesize\textsc{Keywords} --- \textit{#1}}\par}

% --- L'appareil critique ----------------------------------------------------

% Le numéro de feuillet, en marge. C'est l'ancre qui permet de régler un
% désaccord sur une formule en regardant *un* feuillet plutôt qu'une chemise
% de six cents. Le site s'en sert aussi pour tourner les pages du fac-similé
% quand on fait défiler la transcription.
\newcommand{\leaf}[1]{%
  \marginnote{\footnotesize\color{gray!70}\textsf{#1}}%
  \label{leaf:#1}\ignorespaces}

% Illisible : on ne devine pas. Un mot inventé qui se lit comme les autres est
% le pire résultat possible.
\newcommand{\ill}{\textcolor{red!60!black}{[\,\ldots\,]}}

% Lecture douteuse : le mot est donné, et signalé comme incertain.
\newcommand{\uncertain}[1]{\uline{#1}}

% Ajout de l'éditeur — une lettre manquante, un mot sous-entendu.
\newcommand{\add}[1]{\textcolor{blue!50!black}{[#1]}}

% Biffé par Grothendieck lui-même. Ce qu'il a rayé fait partie du document :
% une rature dit souvent où la pensée a changé de direction.
\newcommand{\struck}[1]{\sout{#1}}

% Note du transcripteur, sur l'état matériel du feuillet.
\newcommand{\note}[1]{\par\smallskip{\small\color{gray!60!black}\textit{[#1]}}\par\smallskip}

% Note marginale de Grothendieck, distincte du corps du texte.
\newcommand{\marginal}[1]{\par\smallskip\noindent\hspace{1em}%
  {\small\color{gray!60!black}#1}\par\smallskip}

% --- Titre ------------------------------------------------------------------

\AtBeginDocument{%
  \begin{center}
    {\large\bfseries Fonds Alexandre Grothendieck --- cote n\textsuperscript{o}~\grothFolder}\\[2pt]
    {\itshape\grothFoldertitle}\\[4pt]
    {\small feuillets \grothFirst--\grothLast\ (lot \grothBatch)}\\[2pt]
    {\footnotesize\color{gray!60!black}Transcription établie d'après le fac-similé numérisé
     par l'Université de Montpellier.\\
     Les lectures incertaines sont soulignées, les illisibles notées [\,\ldots\,],
     les ajouts entre crochets.}
  \end{center}
  \vspace{1em}}

\endinput


\folder{23}
\batch{1}
\pages{1}{20}
\dating{1966-1968}
\watermark{Édition de démonstration}
\foldertitle{EGA IV. Compléments : notes manuscrites (s.d.), lettres (1966, 1968), tapuscrit (s.d.).}

\begin{document}

\section{EGA IV, Calcul diff. (compléments)}
\note{les mots du titre sont de sa main, en haut à droite du feuillet 1, une
chemise par ailleurs vierge qui couvre l'ensemble de la cote ; le feuillet ne
reçoit pas de numéro de page}

\subsection{Discriminant etc}
\note{titre de sa main, en haut à droite du feuillet 2, chemise vierge qui
couvre les trois feuillets suivants}

\page{3}
$E$ faisceau loc. libre de rang $r+1$ sur $S$, $P = P(E)$ fibré projectif
associé, on a une suite exacte canonique
\[
  0 \longrightarrow \Omega^{1}_{P/S} \longrightarrow g^{*}(E)(-1)
  \longrightarrow \mathcal{O}_{P} \longrightarrow 0
\]
[qui donne un isom. can.
\[
  \tag{1}
  \Omega^{r}_{P/S} \xrightarrow{\;\sim\;}
  g^{*}\bigl(\Lambda^{r+1}_{\mathcal{O}_S} E\bigr)(-r-1)
\]
]
\note{le morphisme $g : P \to S$ n'est pas nommé sur le feuillet ; le $g^{*}$
des formules le suppose}

Si donc $X$ est un sous-préschéma de $P$ \uncertain{lisse} sur $S$, de dim
relative $n = r - p$, associé au faisceau \struck{\ill{}} d'idéaux $J$, on a
la suite exacte
\note{« $n = r-p$ » est ajouté en marge droite, entre parenthèses, à hauteur
de « dim relative » ; « d'idéaux » est écrit en interligne au-dessus d'un mot
biffé}
\[
  0 \longrightarrow J/J^{2} \longrightarrow \Omega^{1}_{P/S} \otimes \mathcal{O}_X
  \longrightarrow \Omega^{1}_{X/S} \longrightarrow 0
\]
d'où l'isomorphisme
\[
  \tag{2}
  \Omega^{r}_{P/S} \otimes \mathcal{O}_X \simeq
  \Omega^{n}_{X/S} \otimes \Lambda^{p}_{\mathcal{O}_X}(J/J^{2})
\]
\note{l'exposant de $\Omega_{X/S}$ est repassé plusieurs fois, ici et dans
la formule suivante ; on lit $n$, ce que la formule du feuillet 4,
$\Omega^{n}_{X/S}$, confirme}
d'où en composant (1) et (2), désignant par $f : X \to S$ le morphisme
structural
\[
  \struck{\Omega^{r}}\; f^{*}\bigl(\Lambda^{r+1}_{\mathcal{O}_S} E\bigr)(-r-1)
  \simeq \Omega^{n}_{X/S} \otimes \Lambda^{p}_{\mathcal{O}_X}(J/J^{2})
\]

Supposons que $X$ soit une intersection complète, définie par un sous-module
gradué $M \subset$ \add{$\mathrm{Sym}^{*}(E)$} (loc. facteur direct de
$\mathrm{Sym}^{*}(\underline{E})$), \struck{\ill{}} loc. libre de rang $p$
[il n'est \struck{d'ailleurs} déterminé que modulo le carré $\underline{I}^{2}$
de l'Idéal qu'il définit ; mais sa structure est \uncertain{commune}].
\ill{} un hom., si $I_d$ est la composante de degré $d$ de $I$
\note{« gradué $M \subset$ » est en interligne ; le module dont $M$ est
sous-module n'est pas répété après $\subset$, l'ajout est de l'éditeur}
\[
  g^{*}(I_d)(-d) \longrightarrow J
\]
d'où des homs surjectifs
\[
  \coprod g^{*}(M_d)(-d) \longrightarrow J
\]
induisant un isom
\[
  \coprod f^{*}(M_d)(-d) \xrightarrow{\;\sim\;} J/J^{2} \simeq
  \mathrm{Conormal}_{P/X}
\]
\note{l'indice de Conormal est bien écrit $P/X$ sur le feuillet}

\page{4}
Par suite, on a
\[
  \Lambda^{p}_{\mathcal{O}_X}\, J/J^{2} \simeq
  f^{*}\Bigl(\bigotimes_{i} \Lambda^{\alpha_i} M_{d_i}\Bigr)
  \bigl(-\textstyle\sum \alpha_i d_i\bigr)
\]
où $M_{d_i}$ sont les composantes \uncertain{non nulles} de $M$, supposées de
degré $d_i$. D'où un isom
\[
  \Omega^{n}_{X/S} \otimes f^{*}\Bigl(\bigotimes_{i} \Lambda^{\alpha_i}
  M_{d_i}\Bigr)\bigl(-\textstyle\sum \alpha_i d_i\bigr)
  \xrightarrow{\;\sim\;} f^{*}\bigl(\Lambda^{r+1}_{S} E\bigr)(-r-1)
\]
[Correspondant : un hom qui existe en tous cas dès qu'on se donne des
$M_{d_i} \to \mathrm{Sym}^{d_i}(\underline{E})$, les $M_{d_i}$ loc. libres
de rang $\alpha_i$, et $X$ défini par \uncertain{l'idéal engendré par} les
images des $M_{d_i}$, posant $p = \sum \alpha_i$, $n = r - p$], i.e. on a
\note{« l'idéal engendré par » est en interligne au-dessus de « les images »}
\[
  \Omega^{n}_{X/S} \xrightarrow{\;\sim\;}
  f^{*}\Bigl(\bigl(\Lambda^{r+1}_{S} E\bigr) \otimes
  \bigotimes_{i} \Lambda^{\alpha_i} \check{M}_{d_i}\Bigr)
  \bigl(\textstyle\sum \alpha_i d_i - r - 1\bigr)
\]

Dans le cas où
\[
  \textstyle\sum \alpha_i d_i = r+1
\]
on trouve donc
\[
  \Omega^{n}_{X/S} \xrightarrow{\;\sim\;} f^{*}(\ -\ )
\]
\note{la parenthèse est laissée vide sur le feuillet, avec un tiret}

Supposons p. ex. que $E$ soit le faisceau loc. libre de rang $3$ associé
\add{à} une courbe elliptique lisse sur $S$, donc $r = 2$, il n'y a qu'un
seul $M_{d_i}$, $d_i = 3$, $\alpha_i = 1$, on trouve
\[
  \Omega^{1}_{X/S} \xrightarrow{\;\sim\;} f^{*}\bigl(\Lambda^{3} E \otimes
  \check{M}\bigr)
\]
Or $E$ a une filtration de facteurs $\mathcal{O}_S$, \struck{$L$}
$L^{\otimes 2}$, $L^{\otimes 3}$, où $L$ est inversible tq
$f^{*}(L) \simeq \mathcal{T}_{X/S}$, donc $\Lambda^{3} E \simeq L^{\otimes 5}$,
et comme $\Omega^{1}_{X/S} \simeq f^{*}(\check{L})$, on trouve
$\check{L} \simeq L^{\otimes 5} \otimes \check{M}$ d'où $M \simeq
L^{\otimes 6}$.
\note{le symbole après « $f^{*}(L) \simeq$ » est une lettre cursive à boucle,
lue $\mathcal{T}$ (le faisceau tangent) ; la ligne suivante,
$\Omega^{1}_{X/S} \simeq f^{*}(\check{L})$, est cohérente avec cette lecture}

\page{5}
\note{feuillet de brouillon : des formules alignées à gauche, un bloc de
prose au centre, une insertion écrite en diagonale dans la marge gauche ; on
suit l'ordre de lecture du haut vers le bas}

$E$ \qquad $\mathrm{Symm}_k(\underline{E})^{\vee}$

$P^{r}$ \qquad $\underline{\Phi}^{k}(P^{r})$
\note{le $\Phi$ est barré d'un trait horizontal, son trait constant dans ce
feuillet ; on le rend $\underline{\Phi}$}

\[
  \delta_{r,k} \in \underline{\Phi}^{(k-1)(r+1)}\bigl(\Phi^{k}(P^{r})\bigr)
\]
(discriminant) défini au signe près
\marginal{$\underbrace{k-1 \quad k-1 \;\cdots\; k-1}_{r+1 \text{ facteurs}}$
\qquad $(k-1)(r+1)$}
\note{la note marginale est à droite de la formule ; l'accolade porte
« $r+1$ facteurs » et le produit qu'elle annonce est ensuite écrit
$(k-1)(r+1)$, tel quel, ici et dans les deux exposants suivants}

section de $\underline{\mathcal{O}}_{\Phi^{k}(P^{r})}\bigl((k-1)(r+1)\bigr)$

Considérons une section $\varphi$ de $\Phi^{k}(P^{r})$ sur $S$, elle définit
donc un \struck{homom. \ill{} $\Phi^{k}(P^{r})$ section de} faisceau
inversible $\check{M}$ sur $S$, image inverse de
$\underline{\mathcal{O}}_{\Phi^{k}(P^{r})}(1)$] et si \struck{\ill{}}
\struck{l'image \ill{} la section} \struck{\ill{}} une section $\delta$ de
$\check{M}^{(k-1)(r+1)}$.
\marginal{quotient \uncertain{inversible} de $\mathrm{Symm}_k(E)^{\vee}$,
\ill{} sous-faisceau $M$ \ill{} de $\mathrm{Symm}_k(E)$}
\note{l'insertion marginale est écrite en diagonale dans la marge gauche et
reliée par une flèche à « faisceau inversible $\check{M}$ » ; « $\varphi$ »
est ajouté en interligne au-dessus de « section », et le « $\delta$ » de
« une section $\delta$ » de même ; le crochet fermant n'a pas d'ouvrant
visible}
Donc la section $\varphi$ \ill{} par $\Phi^{k}(P^{r})_{\delta_{r,k}}$, i.e.
$P^{r}_{\varphi}$ \uncertain{est} lisse sur $S$, sss $\delta$ est partout
non nulle. Notons que d'ailleurs, si $J = J_{\varphi}$ est l'idéal défini
par $\varphi$ sur $P^{r}$, alors
\[
  J \simeq g^{*}(M)(-k)
\]
\note{au-dessus du $M$ un signe est noirci ; la formule encadrée qui suit,
$\mathcal{O}(-k) \otimes g^{*}(M) \to \mathcal{O}_{P^r}$, donne $M$ sans
dual}
\[
  \Bigl[\ \struck{\mathcal{O}(-k) \otimes g^{*}(E)} \longrightarrow
  \mathcal{O}(-k) \otimes g^{*}\mathrm{Symm}_k(E) \longrightarrow
  \mathcal{O}_{P^{r}} \qquad
  \text{d'où} \quad \mathcal{O}(-k) \otimes g^{*}(M) \longrightarrow
  \mathcal{O}_{P^{r}} \quad -\ \Bigr]
\]
\note{sous le terme biffé, « pour \ill{} d'une » ; en bas à gauche du
feuillet, un griffonnage et deux mots, « \ill{} \ill{} », qui ne se lisent
pas}

\subsection{Exponentielle}
\note{titre de sa main, en haut à droite du feuillet 7, chemise vierge qui
couvre les feuillets 8 à 16}

\page{8}
\textbf{Sorites sur l'exponentielle}

(associative avec unité) \uncertain{linéairement topologisée}, sép. et
complète
\note{la parenthèse et ce qui suit sont écrits en petit sous le titre, à
droite, et raccordés par un trait à « algèbre » de la ligne suivante : ils
qualifient l'algèbre $U$}

(1) Soit $U$ une algèbre [sur $\mathbf{Q}$], $N(U)$ l'ens. de ses éléments
top. nilpotents. On définit deux \emph{bijections} \emph{réciproques} l'une
de l'autre
\[
  N(U) \;\underset{\log}{\overset{\exp}{\rightleftarrows}}\; 1_U + N(U)
\]
par
\[
  \begin{cases}
    \exp \xi = \sum_{n \geqslant 0} \frac{1}{n!}\,\xi^{n} \\[4pt]
    \log(1+\eta) = \sum_{n > 0} (-1)^{n+1} \frac{1}{n}\,\eta^{n}
  \end{cases}
\]
[N.B. Si la top. de $U$ est discrète, $N$ est l'ens. des éléments
nilpotents]. On a
\[
  \begin{cases}
    \exp(\xi + \xi') = \exp \xi \, \exp \xi' \quad \text{si } [\xi, \xi'] = 0 \\[4pt]
    \struck{\log((1+\eta)(1+\eta'))} \\[2pt]
    \log g g' = \log g \; \log g' \quad \text{si } g g' = g' g
  \end{cases}
\]
\note{les deux conditions « si » sont reliées par une double flèche
verticale $\Updownarrow$ ; le second membre de la dernière formule est écrit
tel quel, $\log g \; \log g'$, sans signe $+$ visible}

[Si \struck{\ill{}} $I$ est un idéal (à gauche ou à droite) de $U$,
\uncertain{fermé}, \ill{}
\[
  N(U) \cap I \;\rightleftarrows\; 1_U + N(U) \cap I
\]
Si $\varphi : U \to V$ est un homom. d'\struck{anneaux} \add{algèbres}
topologiques, on a $\varphi(N(U)) \subset N(V)$, et
\[
  \begin{cases}
    \varphi \exp \xi = \exp \varphi(\xi) \\[2pt]
    \varphi \log g = \log \varphi(g)
  \end{cases}
\]
\note{le mot de remplacement au-dessus de « anneaux » ne se lit pas ;
« algèbres » est une conjecture de l'éditeur ; « fermé » est écrit en
interligne au-dessus de « (à gauche ou à droite) »}

\emph{Cas} \uncertain{Toute} sous-algèbre fermée \ill{} est \ill{}
« stable » pour $\exp$, $\log$
\note{le bloc entre crochets qui précède, du « Si $I$ est un idéal » à la
formule $\varphi \log g = \log \varphi(g)$, est accolé d'un grand crochet
dont une flèche descend jusqu'à cette dernière ligne : il est à reporter
après elle}

\page{9}
(2) Supposons que $U$ soit \struck{\ill{}} une \uncertain{k}-\ill{}
augmentée [\uncertain{$k$} \ill{} \ill{} \ill{} \ill{}] \ill{} \ill{}
\uncertain{muni} \uncertain{d'une} \ill{} diagonale
\note{la lettre cerclée en interligne au-dessus du mot biffé est un $k$,
répété une ligne plus bas ; on lit « augmentée » sans certitude}
\[
  \Delta : U \longrightarrow U \,\widehat{\otimes}_{k}\, U
\]
qui soit un \uncertain{hom.} de k-algèbres \uncertain{augmentées}
\uncertain{et} un \uncertain{hom.} \ill{} k-\uncertain{bigèbre}
\struck{\ill{}} (\ill{} \ill{} \ill{}). On \uncertain{désigne} par
\struck{$\mathfrak{g}(U)$} $\gamma(U)$ \struck{$\pi(U)$} l'ens. des
$g \in 1 + I(U)$ \struck{\ill{}} \uncertain{tels} que
\note{trois notations sont écrites l'une sous l'autre pour l'ensemble qui
va être défini : $\mathfrak{g}(U)$, biffé, $\gamma(U)$, et une troisième
biffée par-dessus un $\pi$ souligné ; c'est $\gamma(U)$ qui est repris
aux feuillets 10 à 12 ; « depuis quelques » et une flèche en interligne
au-dessus de « (...) » ne se rattachent à rien de lisible}
\[
  \Delta(g) = g \otimes g
\]
\[
  \delta \xi = \xi \otimes \xi \qquad \text{si} \quad g = 1 + \xi
  \quad (\xi \in I(U))
\]
\note{la formule $\delta\xi = \xi\otimes\xi$ est écrite sous
$\Delta(g) = g\otimes g$ ; la lettre $\delta$ paraît désigner la
« diagonale réduite » sur $I(U)$, qui n'est pas définie sur le feuillet}
\uncertain{correspondants} \ill{} \uncertain{définit} \ill{}
[\struck{\ill{}} \uncertain{ils} \uncertain{correspondent}
\uncertain{bien} à la bigèbre \struck{\ill{}} $k[T]$, avec
$\delta T = T \otimes T$, \uncertain{ou} \uncertain{plutôt} \ill{}
\uncertain{première} $k[[T]]$, avec $\delta T = T \otimes T$, \ill{}
\ill{} \ill{} \ill{} \ill{} \ill{}~$\}$ ; \ill{} \uncertain{le} \ill{}
\uncertain{dual} \ill{} \uncertain{et} \ill{} \ill{} $\mathbb{Z}$ ... ].
\uncertain{Ils} \ill{} \ill{} \ill{} \ill{}
\note{le crochet ouvert avant « ils correspondent » se ferme après
« $\mathbb{Z}$ ... » ; entre les deux, la prose ne se laisse lire que
par bribes}
\marginal{\ill{} $I(U)$ \ill{} $\to 0$ \ill{} $\to \infty$ \ill{}
\struck{\ill{}}}
\note{note marginale à gauche de la ligne « ils correspondent », quatre
lignes courtes dont les formules seules se lisent}
\[
  \Delta(g g') = \Delta(g)\,\Delta(g') = (g \otimes g)(g' \otimes g')
  = g g' \otimes g g'
\]
\[
  \Delta(g^{-1}) = \struck{\ill{}}\,(\Delta(g))^{-1}
  = (g \otimes g)^{-1} = g^{-1} \otimes g^{-1}
\]
\uncertain{de} $\gamma(U)$, et $\gamma(U)$ \uncertain{fonctoriel}
\struck{\ill{}} \uncertain{en} $(k, U) \longrightarrow (k', U')$.
\note{les deux dernières lignes : une phrase dont le sujet est en haut de
la ligne précédente et ne se lit pas ; « Fonctorialité » se lit en
marge sous la formule pour $\Delta(g^{-1})$}

\page{10}
Supposons \struck{\ill{}} que $U$ soit une $\mathbf{Q}$-algèbre.
\struck{Alors} Soit $\pi(U)$ l'ens. des \uncertain{éléments}
\uncertain{primitifs} de $U$. Alors
\note{« $\pi(U)$ » est écrit par-dessus « $\exp\xi$ » biffé ; à droite,
un « Alors » biffé}
\[
  \pi(U) \;\underset{\log}{\overset{\exp}{\rightleftarrows}}\; \gamma(U)
\]
\note{au-dessus de la flèche, un premier essai « $\pi(U) \to \gamma$ »
est biffé}
\uncertain{sont} des \emph{bijections}, inverses l'une de l'autre.

[N.B. il suffit de vérifier que $\exp$ applique $\pi(U)$ dans
$\gamma(U)$, et $\log$ applique $\gamma(U)$ dans $\pi(U)$]. Pour
\uncertain{transporter} \uncertain{les} \uncertain{structures},
\uncertain{on} \uncertain{obtient} \ill{} $\pi(U)$ une loi de
\uncertain{gr.} \ill{} \struck{\ill{}} \ill{} \uncertain{donnée} à
l'aide de \struck{\ill{}} \uncertain{crochets}, \uncertain{par} la
\uncertain{formule} \uncertain{universelle} [\uncertain{raison} \ill{}
\ill{} \struck{\ill{}} \uncertain{de} \uncertain{Campbell-}
\uncertain{Hausdorff} (i.e.\ \ill{} \struck{\ill{}} \ill{}
\uncertain{$k$-bigèbre} \struck{\ill{}} \uncertain{libre} \ill{}
\uncertain{sur} \ill{} \uncertain{d'un} \uncertain{ensemble}
\uncertain{complété}), \ill{} \ill{} \uncertain{définie} \ill{}
$\Lambda$ une k-algèbre infinitésimale.
\marginal{\struck{Hausdorff-Campbell} \uncertain{et} \uncertain{explicitée}
des \uncertain{bijections} \uncertain{réciproques.}}
\note{le bas du feuillet est un enchevêtrement de deux ou trois essais
biffés et de renvois fléchés ; on transcrit dans l'ordre de lecture le
plus vraisemblable et on marque l'incertitude mot par mot. Le nom du
théorème est lu « Campbell-Hausdorff » dans le texte et
« Hausdorff-Campbell » dans la marge (biffé), puis « Hausdorff-Campbell »
aux feuillets 14 et 15}

\uncertain{(3)} Soit \struck{$U$ une k-bigèbre} \uncertain{hyp.} \ill{})
\uncertain{Alors}
\[
  \boxed{\ I(U) \ \ill{} \ \exp\bigl\{\ \ill{} \ \bigr\}\ }
\]
\note{un bloc encadré d'un trait ondulé, de deux lignes, dont seuls
« $I(U)$ » et « $\exp\{$ » se lisent ; il est relié par une flèche à
la ligne suivante}
$\Lambda$ une k-algèbre \uncertain{infinitésimale}. Alors
$\struck{\pi} \otimes_{k} U \otimes_{k} \Lambda$ \ill{} une
\uncertain{k-hyperalgèbre}, \struck{\ill{} \ill{} \ill{}} \ill{}
\uncertain{et} \ill{} \uncertain{augmentée}, \ill{}
\uncertain{d'augmentation} $I(U) \otimes_{k} \Lambda$ \ill{} \ill{}
\uncertain{mais} \uncertain{non} \ill{}
\note{« hyperalgèbre » est une lecture douteuse d'un mot qui commence par
« hyp » ; les dernières lignes du feuillet sont cernées d'un trait et
reliées par une flèche au haut du feuillet 11}

\page{11}
$\struck{\pi_{0}(U)} \struck{\otimes_{k} I(\Lambda)}$ \ill{}
\note{un premier essai, biffé, en tête du feuillet}
\[
  \Bigl[\ \struck{I(U) \otimes_{k} I(\Lambda)} \Bigr] \qquad
  \boxed{U \otimes_{k} I(\Lambda)} \longrightarrow
  1 + I(U) \otimes_{k}
\]
\note{deux lignes reliées par une accolade et prolongées par une flèche
vers « $1 + I(U)\otimes_k$ », dont le second facteur n'est pas écrit}

\uncertain{Alors} \qquad $U \otimes \underline{I}(\Lambda)$ \ill{}
\uncertain{de} \uncertain{puiss.} \uncertain{nilpotente}
\note{le $I$ est ici barré d'un trait horizontal, comme le $\Phi$ du
feuillet 5}

\uncertain{aussi}, \uncertain{d'où}
\[
  U \otimes I(\Lambda)
  \;\underset{\log}{\overset{\exp}{\rightleftarrows}}\;
  1 + U \otimes I(\Lambda)
\]
\struck{\ill{}} \uncertain{bijections} \uncertain{d'une}

[\uncertain{et} \uncertain{si} $I(U) \to$ \ill{}, \ill{} $U$ \ill{}
\[
  I(U) \otimes I(\Lambda)
  \;\underset{\log}{\overset{\exp}{\rightleftarrows}}\;
  1 + I(U) \otimes I(\Lambda)
\ \Bigr]
\]

Supposons \uncertain{maintenant} \uncertain{que} $U$ \uncertain{soit}
\uncertain{une} \uncertain{hyperalgèbre} \ill{} \ill{} \ill{}
\uncertain{primitifs} \uncertain{définit} \ill{} \uncertain{Alors} la
\uncertain{comparaison} \ill{} \ill{} \uncertain{les} \uncertain{primitifs}
\uncertain{de} \ill{} \ill{} \uncertain{comm.} \uncertain{bi-\ill{}} \ill{}
$(U \otimes I(\Lambda)) \cap I(U \otimes \Lambda)$
\struck{$= I(U \otimes \Lambda) \otimes I(\Lambda)$}
\struck{\ill{}} \ill{}
\note{la formule est récrite trois fois par-dessus des ratures ; on ne
garde que le terme de gauche, seul à se lire}
\uncertain{primitifs} \uncertain{sont} \uncertain{dans}
$U \otimes_{k} \Lambda$, \uncertain{et} \ill{} \uncertain{la}
\uncertain{partie} \uncertain{de} \quad
$1 + I(U) \otimes I(\Lambda)$ \struck{$= 1 + \gamma(U \otimes \Lambda)$}
\uncertain{formée} \uncertain{des} \uncertain{éléments}
\uncertain{multiplicatifs}. \uncertain{On} \uncertain{voit}
\uncertain{que} \ill{} \uncertain{les} \uncertain{primitifs} \uncertain{de}
$U \otimes_{k} \Lambda$ qui \ill{} \uncertain{dans} \ill{}
\uncertain{$U \otimes_{k} \Lambda \to U$}, \uncertain{et} \uncertain{les}
\uncertain{éléments} \uncertain{multiplicatifs} \uncertain{de}
$U \otimes_{k} \Lambda$ qui \uncertain{donnent} $1$ \ill{}
\uncertain{la} \ill{} \ill{}
\note{les six dernières lignes sont d'une écriture très rapide, d'une
encre qui s'affaiblit vers le bas ; on transcrit ce qui se laisse deviner
mot par mot, en lecture douteuse, et le reste comme illisible}

\page{12}
\uncertain{6)} \ill{} obtient donc \ill{} des bijections, inverses l'une de
l'autre
\[
  \pi_{0}(U \otimes \Lambda)
  \;\underset{\log}{\overset{\exp}{\rightleftarrows}}\;
  \gamma_{0}(U \otimes_{k} \Lambda)
\]
\note{les lettres $\pi$ et $\gamma$ désignent, d'après les feuillets 11 et
12, les éléments primitifs et les éléments « multiplicatifs » ; l'indice
$0$ n'est pas explicité}

\ill{} \ill{} \uncertain{indiquent} (\ill{} \ill{} \ill{} \ill{}
\uncertain{primitifs} resp. multiplicatifs \ill{} \ill{}
$U \otimes \Lambda \longrightarrow U$. \ill{} \ill{} limite projective
\uncertain{$V$}, \ill{} déduit \ill{} \ill{} \ill{} si $\Lambda$ est une
\ill{} \ill{} \ill{} k-algèbre infinitésimale \ill{} \ill{} d'ailleurs \ill{}
$\Lambda$ \uncertain{est} \uncertain{plus}, \struck{\ill{}} \ill{} \ill{}
$\Lambda$ \ill{} \ill{} \ill{}
\note{six lignes de prose entre les deux formules, dont seuls quelques mots
se lisent ; un trait horizontal barre la fin de la cinquième}
\[
  \pi(U) \otimes_{k} I(\Lambda) \xrightarrow{\;\alpha\;}
  \pi_{0}(U \otimes_{k} \Lambda)
\]
\ill{} \ill{} si $\Lambda$ \ill{} k-plat,
\[
  0 \longrightarrow \pi(U) \longrightarrow \struck{\ill{}}
  \uncertain{I(U)} \xrightarrow{\;\delta\;} I(U) \otimes I(U)
\]
\note{le terme du milieu est repassé et biffé, puis récrit par-dessus ; on
lit $I(U)$ sans certitude}
$U$ \ill{} \uncertain{exact}. \ill{}
\note{la ligne est soulignée deux fois}

donc un \ill{} \ill{} \ill{}
\[
  \pi(U) \otimes_{k} I(\Lambda) \xrightarrow{\;\exp\circ\alpha\;}
  1 + \pi_{0}(U \otimes_{k} \Lambda)
\]
\note{le second membre est écrit $1 + \pi_{0}(\ldots)$, tel quel}
\ill{} isom. \ill{} dans le cas \ill{} \ill{} \ill{} précisions, $U$ une
\uncertain{bigèbre} \ill{} \ill{} : \ill{} le corps, \uncertain{th. de
Cartier} \ill{} \ill{} \ill{} et \uncertain{esp.}) \struck{\ill{}} \ill{}
\ill{} $\pi_{0}(U \otimes_{k} \Lambda)$.

\marginal{$\rho$ [i.e.\ $\Delta$ est \ill{} et assoc., \struck{\ill{}}
$I(U) = \bigcup U_i$, \ill{} $\delta U_i \subset U \otimes U_i$ et
$\gamma(U) = \{e\}$] —
$U$ \ill{} d'un \uncertain{point fixe}, notion \ill{}
$\mathfrak{g} = \pi(U)$ \ill{} —
$\rho'$ les \ill{} d'autre part et \ill{} \ill{} \uncertain{Poincaré-Witt}
\ill{} : $\mathfrak{g} \simeq \pi(U(\mathfrak{g}))$ \ill{} —
(\ill{} \ill{} \ill{}), donc \ill{} voit que $U \simeq U(\mathfrak{g})$
\ill{} $\mathfrak{g} = \pi(U)$ : \uncertain{th. de Cartier} \ill{}
\uncertain{En particulier} —
si $\mathfrak{g}$ est de dim. finie, \ill{} donc \ill{} \ill{}
\uncertain{Cor.} \ill{} \ill{} \ill{} \ill{} \ill{}}
\note{note de sa main dans la marge gauche, écrite de bas en haut sur toute
la hauteur du feuillet, en deux ou trois colonnes serrées ; on n'en restitue
que les formules et les mots sûrs, dans un ordre de lecture incertain.
Le $\mathfrak{g}$ est une lettre gothique bouclée, distincte du $\gamma$ de
$\gamma(U)$}

\page{13}
(4) Supposons que $\mathfrak{g} = L_k(M)$, i.e.\ \ill{} un \ill{}, $U$
\ill{} \ill{} \ill{} \ill{} top. discrète.

Soit $\xi$ \uncertain{nilpotent} de \ill{} $U$, $g = \exp \xi$.
\[
  \xi . a = 0 \;\Longleftrightarrow\; g a = a
\]
\[
  \xi \ \ill{} \ \Longleftrightarrow\ g \ \ill{}
\]
\note{la seconde équivalence porte sur un mot de part et d'autre qui ne se
lit pas}
\[
  M' \ \uncertain{\text{stable}} \ \ill{} \ \xi
  \;\Longleftrightarrow\;
  M' \ \uncertain{\text{stable}} \ \ill{} \ g \qquad (M' \subset M)
\]
\note{les trois équivalences sont accolées d'un crochet à gauche ; une
flèche coudée relie la deuxième à la troisième}

Supposons que \ill{} $M$ soit muni \ill{} \ill{} \ill{} \ill{}
$(x, y) \longmapsto x . y$. Alors
\[
  \xi(x . y) = (\xi x) y + x (\xi y)
  \;\Longleftrightarrow\;
  g(x . y) = (g x)(g y)
\]
\uncertain{dérivations} \qquad\qquad \uncertain{automorphismes} ;
\uncertain{multiplicatives} \qquad\qquad \uncertain{invariants}
\note{deux paires de mots, soulignés, sous les deux membres de
l'équivalence ; la lecture des quatre est douteuse}

\ill{} \ill{} \ill{} dans \ill{} \ill{} \ill{} d'augmentation \ill{}
\ill{} \uncertain{algèbre} contient \ill{} \ill{} \ill{} \uncertain{$k$}
\ill{} \ill{} [donc $\Lambda$ contient $\mathbf{Q}$], $\mathbf{Q}$ \ill{}
soit $M$ un \ill{} \ill{}. Alors \ill{} $u \in L_\Lambda(M)$ \ill{}
$u M \subset I(\Lambda) M$ \ill{} $u$ \uncertain{nilpotent} \ill{}
$M \otimes_{\Lambda} \uncertain{k}$ (de \uncertain{puiss.} \ill{} nulle)
\ill{} \ill{} $u \in L_{\Lambda}^{\uncertain{*}}(M)$ \ill{} \ill{}
\ill{} $M \otimes_{\Lambda} k$ \ill{} \ill{} \ill{} \ill{} l'identité
\ill{} \ill{} \ill{}, si \ill{} \ill{} \ill{} \ill{} $\pi_{\ill{}}$
\uncertain{resp.} $\gamma$,
\note{le bas du feuillet est un brouillon dont la prose de liaison ne
passe pas ; « donc $\Lambda$ contient $\mathbf{Q}$ » est relié par un trait
au « $\mathbf{Q}$ » de la ligne suivante, et « $M \otimes_{\Lambda} k$ » est
en interligne au-dessus de « (de puiss. ... nulle) »}

\page{14}
\ill{} obtient des bijections, inverses l'une de l'autre,
\[
  \pi \;\underset{\log}{\overset{\exp}{\rightleftarrows}}\; \gamma
  \qquad (\uncertain{\text{multiplicatif}})
\]
Bien entendu, $\pi$ est un idéal (de $L_{\Lambda}(M)$), et
\uncertain{à fortiori} une \uncertain{sous-}algèbre de Lie. En
\uncertain{transportant} à $\pi$ la structure de \uncertain{groupe} de
$\gamma$, \ill{} \uncertain{trouve} une structure de \uncertain{groupe}
dans $\pi$, qui \ill{} \ill{} \ill{} \uncertain{de} \ill{} structure
d'algèbre de Lie par la formule de Hausdorff-Campbell. Si $M$ est une
$\Lambda$-algèbre, \ill{} \uncertain{trouve} \ill{} \uncertain{contient}
les \uncertain{automorphismes} de la $\Lambda$-algèbre $A$ (i.e.\ \ill{}
$\Lambda$ \ill{} une algèbre \ill{} \ill{} : \struck{\ill{}} idéal
d'\uncertain{augmentation} \ill{} \ill{}, \ill{} \ill{} une
$\mathbf{Q}$-algèbre) qui induisent l'identité sur
\struck{\ill{}}$\otimes_{\Lambda} k$, correspondent
\uncertain{biunivoquement} aux $\Lambda$-dérivations $D$ de $B$ telles que
$D B \subset \uncertain{I(\Lambda) B}$.
\note{« (multiplicatif) » est écrit sous le $\gamma$ et relié par un trait
à « idéal » ; un « $A$ » est ajouté au-dessus de « $\Lambda$-algèbre » à sa
première occurrence, et le nom de l'algèbre est ensuite $A$ puis $B$ ;
« $D$ » est ajouté au-dessus de « dérivations » ; « contient » et
« automorphismes » sont soulignés}

\page{15}
Supposons en particulier
\[
  B = A \otimes_{k} \Lambda
\]
\uncertain{où} $\Lambda$ est une k-algèbre infinitésimale. Alors
\struck{\ill{} $\Omega^{\prime}_{\Lambda}$}
\[
  \mathcal{T}_{B/\Lambda} \struck{(B/A)}
  = \operatorname{Hom}_{B}(\Omega^{1}_{B/\Lambda}, B)
  \simeq \operatorname{Hom}_{B}(\Omega^{1}_{A/k} \otimes_{k} \Lambda,
    A \otimes_{k} \Lambda)
  \simeq \operatorname{Hom}_{A}(\Omega^{1}_{A/k}, A \otimes_{k} \Lambda)
  \simeq \text{k-dérivations de } A \text{ dans } A \otimes_{k} \Lambda
\]
\note{la lettre cursive à boucle devant $(B/A)$ est celle du feuillet 4,
lue $\mathcal{T}$ : ici elle est explicitée par
$\operatorname{Hom}_{B}(\Omega^{1}_{B/\Lambda}, B)$, ce qui fixe la
lecture pour tout le lot ; « $(B/A)$ » est biffé et remplacé par l'indice
$B/\Lambda$ ; « de $A$ dans $A\otimes_k\Lambda$ » est écrit en petit sous
« k-dérivations »}
et les $\Lambda$-dérivations de $B$ \uncertain{en} $B$ qui induisent
$0$ sur $B \otimes_{\Lambda} k \simeq \Lambda$ s'identifient aux
k-dérivations $A \longrightarrow A \otimes_{k} I(\Lambda)$.
\note{on lit $\Lambda$ après $\simeq$, où l'on attendrait $A$}

Donc

\textbf{Proposition} Supposons $k \struck{\ill{}} \supset \mathbf{Q}$, $A$
une k-alg. \uncertain{comm.}, $\Lambda$ une k-alg. \ill{}. Alors les
\struck{$\Lambda$-aut.} \uncertain{$\Lambda$-automorphismes}
infinitésimaux de $A \otimes_{k} \Lambda$ \struck{\ill{}} \ill{}
\uncertain{correspondent} à l'\uncertain{ensemble} des k-dérivations de
$A$ dans $A \otimes_{k} I(\Lambda)$, \ill{} la structure de
\uncertain{groupe} donnée par Hausdorff-Campbell.
\note{« $\supset \mathbf{Q}$ » est écrit par-dessus un mot biffé, le
$\mathbf{Q}$ cerclé ; le mot après « $\Lambda$ une k-alg. » est une
abréviation de trois ou quatre lettres qui ne se lit pas ; une insertion
au-dessus de « $\Lambda$-aut. » ne se lit pas non plus. Un trait vertical
ondulé accole tout l'énoncé dans la marge gauche, et se poursuit sur les
trois premières lignes du feuillet 16}

\page{16}
En particulier, si $\Lambda$ est k-libre de \ill{} \ill{} \ill{}
\uncertain{fini}, \ill{} \uncertain{alors}
\[
  \mathcal{T}_{A/k} \otimes_{k} I(\Lambda)
  \simeq \pi_{0}\bigl(\uncertain{\mathcal{T} \otimes_{k} U}\bigr),
\]
\ill{} \ill{} \ill{} \ill{} $=$ \ill{} si $\Omega^{1}_{A/k}$ est libre de
\ill{} \ill{} fini sur $A$.
\note{le second membre $\pi_{0}(\ldots)$ est écrit en interligne
au-dessus, relié par une accolade à la formule ; la ligne « $= \ldots$ »
qui suit ne se lit pas}

[alors $A/k$ est \uncertain{diff.} \uncertain{simple} \ill{}]

\subsection{Fibrés des germes d'ordre $n$ de sections, etc. Transformations
infinitésimales}
\note{titre de sa main, en haut à droite du feuillet 17, chemise vierge qui
couvre les feuillets 18 à 20 et sans doute la suite du lot 2 ; les mots
« sections » et « infinitésimales » sont soulignés}

\page{18}
$S$ préschéma

$\Lambda$ faisceau loc. libre de rang fini de
$\underline{\mathcal{O}}_{S}$-algèbres
$\underline{\mathcal{O}}_{S}$-augmentées, \struck{tel que} d'idéal
d'augmentation $I$, tel que $\operatorname{gr}^{I}(\Lambda)$ soit
\ill{} localement libre [donc $I^{n} = 0$ pour $n$ \uncertain{grand}].
Un tel $\Lambda$ est dit « faisceau de type infinitésimal ».

$T = \operatorname{Spec}(\Lambda)$, S-schéma fini et plat, muni d'une
section $S \to T$.

\begin{tikzcd}
  X \arrow[d, "f"'] & X_{0} \arrow[l] \arrow[d, "f_{0}"] \\
  S & T \arrow[l] & S = T_{0} \arrow[l]
\end{tikzcd}

\note{le diagramme est encadré dans la marge gauche ; sur la ligne du bas
la flèche de $S = T_0$ vers $T$ est épaissie ; une inscription en biais le
long du cadre, « \ill{} $X_0$ », relie le cadre au « $X$ » du texte et
ne se lit pas}

\[
  \Lambda_{n} = \Lambda / I^{n+1}, \qquad
  T_{n} = \operatorname{Spec}(\Lambda_{n}) \quad \text{donc}
\]
\[
  T_{0} \simeq S \subset T_{1} \subset T_{2} \cdots \subset T_{n}
  \subset \cdots T_{N} = T
\]
$X$ un $T$-schéma, \uncertain{plat} \uncertain{simple} sur $T$,
$X_{n} = X \times_{T} T_{n}$ \struck{\ill{}}
\note{l'insertion au-dessus de « sur $T$ » est de deux mots dont le second
seul se lit à peu près}

On considère de plus
\[
  \xi = X_{0} \struck{\times_{S}} T = X_{0} \otimes_{\underline{\mathcal{O}}_S} \Lambda
  \qquad\Big/\quad \text{Donc } \xi_{n} = X_{0} \times_{S} T_{n},
  \ \text{en particulier } \xi_{0} = X_{0}
\]
Du point de vue \ill{} topologique, \ill{}
\[
  |X| = |\xi| = |X_{0}| = |\xi_{0}|
\]
$\underline{\mathcal{O}}_{X}$ est un faisceau (de
$f_{0}^{*}(\underline{\Lambda})$-algèbres) sur \struck{$X$} $X_{0}$,
augmenté vers $\underline{\mathcal{O}}_{X_0}$,

\begin{tikzcd}
  \underline{\mathcal{O}}_{X} \arrow[r] & \underline{\mathcal{O}}_{X_0} \\
  \underline{\mathcal{O}}_{S} \arrow[r] \arrow[u] &
  \underline{\Lambda} \arrow[r, "\varepsilon"] \arrow[u] &
  \underline{\mathcal{O}}_{S}
\end{tikzcd}

avec commutativité

de plus dans la filtration $I$-adique de $\underline{\mathcal{O}}_{X}$,
\uncertain{pour} le gradué associé est loc. libre sur
$\underline{\mathcal{O}}_{X_0}$ ($X$ plat sur $T$).
\note{« $I$-adique » est en interligne ; dans le diagramme le sommet
$\underline{\mathcal{O}}_{S}$ de droite est le but de $\varepsilon$ et n'a
pas de flèche montante}

\uncertain{Si} pour \uncertain{tout} \ill{} $U \subset X_{0} = |X|$, soit
$\underline{\mathrm{Aut}}^{X_0}_{T}(X)$ \ill{}
$\underline{\mathrm{Aut}}^{X_0}_{\Lambda}(X)$ le groupe des
$T$-automorphismes de $U$ induisant l'identité sur
$X_{0} \cap U = U \times_{T} T_{0}$.
\note{les deux $\underline{\mathrm{Aut}}$ sont écrits l'un sous l'autre
avec un tiret devant le second : deux notations pour le même faisceau}

\marginal{\uncertain{mettre} les \uncertain{hyp.} \ill{} \ill{} pour
\uncertain{simplifier} \ill{} \ill{} \ill{} \ill{}}
\note{note de sa main écrite en diagonale dans la marge gauche, sous le
diagramme encadré ; quatre lignes courtes, en grande partie illisibles}

\page{19}
Ce groupe s'identifie aussi au groupe des
$f_{0}^{*}(\underline{\Lambda})$-automorphismes de
$\underline{\mathcal{O}}_{X} | U$, qui induisent l'identité dans
$\underline{\mathcal{O}}_{X} / \underline{I} . \underline{\mathcal{O}}_{X}$ ;
donc, quand $U$ varie, ce groupe forme les sections d'un faisceau de
groupes, noté
\[
  \underline{\mathrm{Aut}}^{X_0}_{T}(X) \ \ill{} \ =
  \underline{\mathrm{Aut}}^{X_0}_{\Lambda}(X)
\]
qui s'identifie aussi au faisceau des groupes des \add{automorphismes de
l'} $f_{0}^{*}(\underline{\Lambda})$-\struck{Algèbre}
$\underline{\mathcal{O}}_{X_0}$-augmentée $\underline{\mathcal{O}}_{X}$.
\note{« automorphismes de l' » est une insertion de sa main en interligne,
restituée entre crochets parce que sa place exacte dans la phrase n'est
marquée que par un trait}

D'autre part, considérons, \struck{\ill{}} le $S$-schéma
\[
  \overline{X} = \underline{\operatorname{Hom}}_{T/S}(T, X)
\]
\marginal{« schéma des sections de $X$ sur $\Lambda$ »}
\note{la note marginale est à gauche de la formule ; l'argument $T$ de
$\underline{\operatorname{Hom}}$ porte une rature ; « le $S$-schéma » est
en interligne au-dessus d'un mot biffé}
\uncertain{L'homom.} section $S \xrightarrow{\;\sim\;} T_{0} \to T$
définit
\[
  \underline{\operatorname{Hom}}_{T/S}(T, X) \longrightarrow
  \underline{\operatorname{Hom}}_{T_0/S}(T_{0}, X \times_{T} T_{0})
  \simeq X_{0},
\]
d'où un homomorphisme de $S$-préschémas
\note{« pré » est ajouté en interligne au-dessus de « schémas »}
\[
  \overline{X} \xrightarrow{\;\pi^{X}_{\Lambda}\;} X_{0}
\]
Si $U$ est un ouvert de $X_{0}$, alors on a
\[
  \struck{U}\ \overline{U} \simeq (\pi^{X}_{\Lambda})^{-1}(U)
\]
(compatibilité avec les \uncertain{projections}). Par suite,
$\underline{\mathrm{Aut}}^{X_0}_{T}(X)$ opère comme faisceau
d'automorphismes sur le \uncertain{$X_0$}-schéma $\overline{X}$

\page{20}
et par suite aussi \struck{\ill{}} sur le faisceau
$\boxed{\underline{S}(\overline{X}/X_{0})}$ des \struck{\ill{}} sections
de $\overline{X}$ sur $X_{0}$.

On peut aussi considérer le faisceau
$\underline{\mathrm{Is}}^{X_0}_{T}(\struck{\xi}\ \xi, X)$, qui est un
faisceau principal homogène (\uncertain{vérifiant} ...) sous
$\underline{\mathrm{Aut}}^{X_0}_{T}(X)$ (opérant à gauche) et sous
$\underline{\mathrm{Aut}}^{X_0}_{T}(\xi)$ (--- droite). D'ailleurs on a
une section privilégiée $s_{0}$ de $\underline{S}(\overline{\xi}/\xi_{0})$
et on en déduit un hom de faisceaux
\[
  u \longmapsto u(s_{0})
\]
\[
  \underline{\mathrm{Is}}^{X_0}_{T}(\xi, X) \longrightarrow
  \underline{S}(\overline{X}/X_{0})
\]
\note{« $s_0$ » est ajouté en interligne au-dessus de « privilégiée »,
avec une parenthèse ouvrante ; le $\xi$ biffé dans
$\underline{\mathrm{Is}}(\xi, X)$ est récrit à côté}

\textbf{Proposition} L'homomorphisme précédent est un isomorphisme de
faisceaux, \uncertain{compatible} avec \struck{\ill{}} \ill{} faisceau de
\ill{} \uncertain{l'opération} (\uncertain{à} gauche) \ill{}
$\underline{\mathrm{Aut}}^{X_0}_{T}(X)$ (à droite).
\note{l'énoncé est accolé d'un trait vertical dans la marge gauche ; la
seconde ligne est en partie biffée et récrite}

\textbf{Corollaire 1} $\underline{S}(\overline{X}/X_{0})$
\uncertain{est} principal homogène sous le faisceau
$\underline{\mathrm{Aut}}^{X_0}_{T}(X)$ (à droite).

\textbf{Corollaire 2} $\underline{\mathrm{Aut}}^{X_0}_{T}(\xi)$ opère de
façon naturelle sur $\underline{S}(\overline{X}/X_{0})$, (qui devient un
faisceau principal hom. \ill{} \ill{} \ill{} faisceau de groupes \ill{}
\ill{} \ill{}).
\note{la dernière ligne du feuillet est coupée par le bord du scan}

\marginal{ceci est \ill{} \ill{} \ill{} \ill{} \ill{} si $\Lambda$ \ill{}
\ill{} $\underline{\mathcal{O}}_{S}$-algèbre \ill{} \ill{} \ill{} \ill{},
\ill{} \ill{} \uncertain{Greenberg} \ill{} \ill{} de $\overline{X}$}
\note{note de sa main dans le coin inférieur gauche, écrite en biais et
cernée d'un trait ondulé, avec une flèche vers le corollaire 2 ; cinq
lignes dont un nom propre, lu Greenberg sans certitude, est le seul mot
qui se laisse deviner}

\end{document}
